\documentclass[output=paper,colorlinks,citecolor=brown,babel,arabicfont]{langscibook}
\ChapterDOI{10.5281/zenodo.18845608}
\author{Javier Leibiusky\orcid{}\affiliation{LOEWE (Goethe Universität) and CERMOM (INALCO)}}
\title[{\textit{Meʿam Loʿez} of \textit{Pirqey ʾAvot} by Itzhak Magriso (Istanbul 1753)}]{{\textit{Meʿam Loʿez} of \textit{Pirqey ʾAvot} by Itzhak Magriso (Istanbul 1753): Scripts, language and identity}}
\abstract{\textit{Meʿam Loʿez} was an influential 18th-century literary project initiated by Yaakov Huli in the Ottoman Empire, aiming to provide a comprehensive Bible commentary in Judeo-Spanish, the language of Ottoman Jews. Huli began with Exodus but passed away before completing it. Itzhak Magriso continued, completing commentaries on Exodus, Leviticus, and Numbers. Magriso also included a commentary on \textit{Pirqey ʾAvot}, a Mishnaic tractate integral to the Jewish diaspora's synagogue readings.

Written in Judeo-Spanish with features of orality, such as repetitions, direct speech, and frequent code-switching between Hebrew, Judeo-Spanish, Ladino (calque Judeo-Spanish), and Turkish, the text uses both Rashi (Sephardic) and Hebrew square scripts, reflecting complex identity questions. A Latin transcription required critical decisions linked to these identity dynamics.

This article explores these identity issues, code-switching, and linguistic borrowing, focusing on how these elements support the author's goal of educating readers with accessible religious and conceptual tools. While Hebrew is deeply embedded -- through loanwords, calques, and rabbinical concepts -- Magriso's commentary, compared to Huli’s, is occasionally less clear pedagogically, as it demands a higher degree of Hebrew literacy.

The work demonstrates Hebrew's prominence within the text, balancing educational objectives with identity affirmation. The author often translates Hebrew terms into Judeo-Spanish for clarity, but when untranslated, context aids comprehension. The unchanged morphology of Hebrew loanwords highlights a strong identity marker, aligning with the text's serious tone and its dual purpose of education and cultural preservation.
}

\IfFileExists{../localcommands.tex}{
 \addbibresource{../localbibliography.bib}
 \usepackage{tabularx,multicol}
\usepackage{url}
\urlstyle{same}
\usepackage{textcomp}
 
\usepackage{langsci-optional}
\usepackage{langsci-lgr}
\usepackage{langsci-gb4e}
\usepackage{langsci-branding}

% \usepackage{fontspec}
\babelprovide{georgian} % Georgian language
\babelprovide{armenian} % Armenian language
\babelfont[georgian,armenian]{rm}{DejaVu Sans}
\newfontfamily\georgianfont{DejaVu Sans}
\newfontfamily\armenianfont{DejaVu Sans}

\usepackage{dialogue}
\usepackage{attrib}
\usepackage{attrib} % speaker labels in Koca's chapter
\usepackage{phoenician} % alf, bet in Mengozzi's chapter

 

\makeatletter
\let\thetitle\@title
\let\theauthor\@author
\makeatother

\newcommand{\togglepaper}[1][0]{
   \bibliography{../localbibliography}
   \papernote{\scriptsize\normalfont
     \theauthor.
     \titleTemp.
In Saloumeh Gholami \& Ulrich Mehlem (eds.), \textit{Identity formation, language and migration dynamics: Minorities in the MENA region and the German diaspora}, Berlin: Language Science Press [preliminary page numbering]
   }
   \pagenumbering{roman}
   \setcounter{chapter}{#1}
   \addtocounter{chapter}{-1}
}

\newbool{bookcompile}
\booltrue{bookcompile}
\newcommand{\bookorchapter}[2]{\ifbool{bookcompile}{#1}{#2}}


\newcommand{\georgian}[1]{{\georgianfont #1}}
\newcommand{\armenian}[1]{{\armenianfont #1}}

\newcommand{\textarab}[1]{{\arabicfont \beginR #1} \endR}
\newcommand{\texthebrew}[1]{{\hebrewfont \beginR #1} \endR}
\newcommand{\textsyriac}[1]{{\syriacfont \beginR #1} \endR}


\newfontfamily\phoenicianfont{Tuffy Regular.ttf}

 %% hyphenation points for line breaks
%% Normally, automatic hyphenation in LaTeX is very good
%% If a word is mis-hyphenated, add it to this file
%%
%% add information to TeX file before \begin{document} with:
%% %% hyphenation points for line breaks
%% Normally, automatic hyphenation in LaTeX is very good
%% If a word is mis-hyphenated, add it to this file
%%
%% add information to TeX file before \begin{document} with:
%% %% hyphenation points for line breaks
%% Normally, automatic hyphenation in LaTeX is very good
%% If a word is mis-hyphenated, add it to this file
%%
%% add information to TeX file before \begin{document} with:
%% \include{localhyphenation}
\hyphenation{
    par-a-digm
    Ver-to-vec
    Go-go-lin
    sa-gen
    mo-misch
    ant-wor-te-te
    dix-wa-zim
    spre-chen
    B-Verf-GE
    Swe-dish
    re-mains
    schrei-ben
    di-pey-vin
    Chris-ten-sen
    Wes-sen-dorf
    Ja-va-khi-shvi-li
    ana-lysis
}

\hyphenation{
    par-a-digm
    Ver-to-vec
    Go-go-lin
    sa-gen
    mo-misch
    ant-wor-te-te
    dix-wa-zim
    spre-chen
    B-Verf-GE
    Swe-dish
    re-mains
    schrei-ben
    di-pey-vin
    Chris-ten-sen
    Wes-sen-dorf
    Ja-va-khi-shvi-li
    ana-lysis
}

\hyphenation{
    par-a-digm
    Ver-to-vec
    Go-go-lin
    sa-gen
    mo-misch
    ant-wor-te-te
    dix-wa-zim
    spre-chen
    B-Verf-GE
    Swe-dish
    re-mains
    schrei-ben
    di-pey-vin
    Chris-ten-sen
    Wes-sen-dorf
    Ja-va-khi-shvi-li
    ana-lysis
}

 \boolfalse{bookcompile}
 \togglepaper[8]%%chapternumber
}{}

\begin{document}
\maketitle

\section{Introduction}\label{ch3.7.ss.1}

\textit{Meʿam Loʿez} was a very important literary project initiated by Yaakov Huli at the beginning of the 18th century in the \ili{Ottoman} Empire. The aim of this project was to make a linear commentary of the Bible in \ili{Judeo-Spanish}, using all Jewish literature available at that time, thus allowing the Jewish community in the \ili{Ottoman} empire to have access to all this literature, which was in \ili{Hebrew}, and therefore inaccessible for most \ili{Ottoman} Jews.
\largerpage[2]

For some scholars, namely Aldina \citet{quintana2006} or Rosa \citet{asenjo2003}, the project was one of several attempts to repair the damage caused by the Sabbetai Zevi affair during the 17th century and the resulting great crisis of Jewish communities. Other scholars, such as Mathias \citet{lehmann2014} or Avigdor \citet{levy1992}, insist that Zevi's affair did not have such an impact on Jewish education, and that the reason for Huli's project was rather a part of a movement that intended to decentralize Jewish education, allowing community members to study by themselves without depending on rabbinical authorities, and this, influenced by Lurianic kabbalistic ideas concerning the importance of individual work for the group.


Huli aimed to produce an educational tool, an extensive commentary of the Bible in the language of \ili{Ottoman} Jews, \ili{Judeo-Spanish}. Huli himself was only able to begin the project, he died while working on the commentary on Exodus. Itzhak Magriso took his place and finished the commentary on Exodus, Leviticus, and Numbers. Into the commentary of Leviticus he embedded a commentary on \textit{PIRQEY ʾAVOT}, a tractate from the \textit{Mishna} (oral tradition) that had become part of the yearly reading cycle in synagogues throughout the Jewish \isi{diaspora}. This commentary was first published in 1753 in Constantinople. Further editions were printed in Livorno -- 1822, Salonica -- 1802 and 1866 and Izmir -- 1866.


The first edition is accessible online through the website of the National Library of Israel (system number 990012377230205171). The commentary of \textit{Pirqey ʾAvot} starts at page 98 and spreads through 55 double pages until page 153 \citep{magriso1753}. 

\begin{figure}
  \centering
  \includegraphics[height=.4\textheight]{figures/ch3.7.fig1.png}
  \caption{First edition cover (Constantinople, 1753)}
  \label{ch3.7.fig.1}
 \end{figure}

\begin{figure}
  \centering
  \includegraphics[height=.4\textheight]{figures/ch3.7.fig2.png}
  \caption{Internal double-page}
  \label{ch3.7.fig.2}
 \end{figure}
 
\textit{Meʿam Loʿez} is written in \ili{Judeo-Spanish} as it was spoken in the community \citep{quintana2005} where the text was written, and it presents different features of orality, such as errors, repetitions, direct speech, and of course, many cases of code-switching between \ili{Hebrew}, \ili{Judeo-Spanish}, \ili{Ladino} (a calque language, \citealt{sephiha1973}), and some \ili{Turkish}.

Generally speaking, Jews in different geographical areas have always chosen to write their Judeo languages in \ili{Hebrew} characters. This was the case for European Jews in the Middle Ages, as well as in the Iberian Peninsula. 

According to Robert L. \citet[128--129]{cooper1989}, in social linguistics, we recognize that the written form of a language has great social and symbolic importance. Even if, theoretically, a group of speakers can choose any writing system, provided it enables them to represent their phonological system sparingly and clearly \citep[112]{bunis2004}, the written signs chosen also function as elements of identity. 

The Jews of medieval Spain had adopted new Romance languages, but they rejected the \ili{Latin} \isi{alphabet}, which for them symbolized the Catholic religion. For this reason, until the beginning of the 20th century, \ili{Judeo-Spanish} speakers kept the \ili{Hebrew} \isi{alphabet} to write their language. 

What is interesting is that after their expulsion from Spain in 1492, the Jews took their vernacular Romance languages with them and, following a process of koinéisation \citep[23]{minervini2006}, not only preserved the language but also gave it the status of a Jewish affective language, the language of the intimacy of the home and the family \citep[83]{bornes-varol2011}.

The fact that they continued to write their language in \ili{Hebrew} characters certainly stems from the mystical character given in Jewish tradition to these letters, as we can see, for example, in the Talmud (Berakhot 55:1):

\begin{quote}
Rav Yehuda said that Rav said: Bezalel knew how to join the letters with which heaven and earth were created. From where do we derive this? It is written here in praise of Bezalel: ``And I have filled him with the spirit of God, in wisdom, and in understanding, and in knowledge, and in all manner of workmanship'' (Exodus 31:3).
\end{quote}

Or in \textit{Sefer Yetzirah} 2:2:
\begin{quote}
Twenty-two letters: He engraves them, carves them, weighs and transmutes them, combines them and forms with them the soul of all created, and the soul of all He has yet to form.
\end{quote}

According to \citet[80]{bornes-varol2011}, the fact that \ili{Judeo-Spanish} people were the only ones to speak the \ili{Judeo-Spanish} language in the \ili{Ottoman} Empire is also an element of identity construction. 

On the other hand, speakers of \ili{Judeo-Spanish} have a lost identity language, the \ili{Hebrew} of religious texts, and this, as we shall see later, is a very important identity feature of the text we are analyzing, materializing in the form of the many code-switching cases, borrowings and calques, which go far beyond the presence of these elements in the spoken vernacular. 

And while we are dealing with the links between writing and identity, a word should be said about the choices I made in transcribing \ili{Judeo-Spanish} and \ili{Hebrew}.

I have adopted a transcription standard for \ili{Judeo-Spanish}, called \textit{akí yerushalayim}. The \textit{akí yerushalayim} spelling was developed in an effort to homogenize the spellings used for \ili{Judeo-Spanish}, with two important criteria \citep{shaul1985}: 

\begin{itemize}
 \item To reflect the phonological system of \ili{Judeo-Spanish} as faithfully as possible, and in particular the many phonemes that do not exist in modern \ili{Spanish}; and 
 \item That the system adopted is easy to learn and assimilate for readers all over the world.
\end{itemize}

This system is easy to reproduce and requires no diacritical or special signs. At the same time, it is the closest to the actual pronunciation of \ili{Judeo-Spanish}. This is important, as pronunciation is a key element of identity for \ili{Judeo-Spanish} people, one of the features that differentiates their language from that of the \ili{Spanish}.

For the \ili{Hebrew}, I have followed the system proposed by the \textit{Revue des  études juives} in France. Again, this is a practical approach whose aim is to make the text as accessible as possible, while maintaining scientific rigor. One of the great advantages of this system is that it uses very few diacritical marks, making reading much simpler and more accessible. 

Non-integrated \ili{Hebrew} loanwords are shown here in italics and small capital letters, while integrated terms are shown only in italics. 

Each example given, which is obviously taken from the corpus, is accompanied by the reference where it appears, so ``1:7'' refers to the commentary on the first chapter, seventh article (\textit{mišnaʰ}) in Magriso's work. 

The article will start by presenting the two different \ili{Hebrew} alphabets used in the text, that will establish the base for reflecting on the relationship between \isi{script} and identity, a question that will be treated later.

The main part of this article follows, which deals with the important presence of \ili{Hebrew} in the text. I will show that \ili{Hebrew} is present not only in loanwords and code-switching cases, but also in a deeper layer, the calque language known also as \textit{Ladino}. The particularity in this case is that \ili{Hebrew} was not really a contact language for \ili{Judeo-Spanish} speakers of the \ili{Ottoman} Empire, since they neither used the language in everyday life, nor really were able to read the traditional texts, except for the rabbis, of course. In this case, we are dealing with a contact language which is only literary, and which is seen as a ``sacred language'', a referential language used to justify and give more prestige to the author's commentaries. One of the best proofs of this is the fact that the author never integrates \ili{Hebrew} loanwords morphologically, contrary to \ili{Turkish} or \ili{Greek} loanwords which are completely integrated. 

Finally, I will present my conclusions on all the topics treated in this article. 

\section{One text, two alphabets}\label{ch3.7.ss.2}

\textit{PIRQEY ʾAVOT}'s commentary in the \textit{MEʿAM LOʿEZ} has the peculiarity of using two different \ili{Hebrew} alphabets. The original text of \textit{PIRQEY ʾAVOT} is in the square \ili{Hebrew} \isi{alphabet}. This not only enhances the text, but also gives the square characters the status of a ``sacred language''.

The commentary, written in \ili{Judeo-Spanish}, is in \ili{Hebrew} Rashi \isi{script}, a form of Sephardic cursive.

According to \citet[129--130]{bunis2004}, up until the time the text was written and published, in the 18th century, \ili{Judeo-Spanish} texts intended for the popular classes were generally written in square \ili{Hebrew} characters, while texts intended for the more learned classes were written in non-vocalized Rashi characters. And it was precisely at this time that the Rashi \isi{alphabet} became the main \isi{alphabet} for \ili{Judeo-Spanish} texts, which helps us to understand why it is used in a text like the \textit{MEʿAM LOʿEZ}, aimed at the popular classes. 

\begin{figure}
  \centering
  \includegraphics[width=\linewidth]{figures/ch3.7.fig3.png}
  \caption{Structure of the text and example of the two scripts}
  \label{ch3.7.fig.3}
 \end{figure}

In \figref{ch3.7.fig.3} we can see a close-up view of the text with a clear representation of the structure of the text and of both scripts used by the author: 

\begin{enumerate}
 \item \ili{Hebrew} printed or square \isi{script} for the original \ili{Hebrew} \textit{Pirqey ʾAvot} articles. 
 \item \ili{Hebrew} Rashi \isi{script}, for the references that the author mentions, in \ili{Hebrew}.
 \item \ili{Hebrew} Rashi \isi{script} for the commentary itself which is mostly in \ili{Judeo-Spanish}.
\end{enumerate}


\ili{Judeo-Spanish} speakers throughout the \ili{Ottoman} Empire were a multi-refer\-enced community in all that concerns identity issues. Beyond the fact of being Jewish -- and the use of the \ili{Hebrew} \isi{script} for writing their language was a clear sign of identity, which we also can find in Judeo-\ili{Arabic}-speaking communities in Northern Africa or in the \ili{Yiddish}-speaking communities in Central and Eastern-Europe -- there was also \textit{Sefarad}, Middle-Age Spain. It is very interesting to note that during the \textit{koinéisation} process, which took place between 16th and 17th centuries, variants of the language depended on the origin of the congregation which held power over the community at the time \citep[20]{minervini2006}. At the beginning, Jews expelled from Spain remained loyal to their cities of origin, but with the time, they mingled and mixed, creating a new Judeo-\ili{Ottoman} identity, which was profoundly anchored in their language and the use of the \ili{Hebrew} \isi{alphabet}. But there was another identity element in the use of the \ili{Hebrew} \isi{alphabet}: it helped \ili{Ottoman} Jews to distinguish themselves from this same Spain they yearned for. Since the basis of their language was Medieval Castilian, the \ili{Judeo-Spanish} of the \ili{Ottoman} Empire distinguished themselves from Spain and \ili{Spanish} through the \ili{Hebrew} \isi{script}, and in so doing, preserved some archaic words and innovations.


It is also interesting to note that almost a century after the commentary of Magriso was published, the \ili{Judeo-Spanish} language went through a new revolution, due to the influence of \ili{French}, through the network of schools implanted by the \textit{Alliance Israélite Universelle} throughout the \ili{Ottoman} Empire. This revolution ended by the adoption of the \ili{Latin} \isi{alphabet} at the beginning of the 20th century. 

\section{The presence of Hebrew}\label{ch3.7.ss.3}

\ili{Hebrew} is present everywhere and at every level in this work. On the one hand we have, of course, the commented text, \textit{PIRQEY ʾAVOT}, in which we find \ili{Biblical Hebrew} in the many references and quotations, and then the text itself, which is written in Mishnaic \ili{Hebrew}. The references Magriso gives in the margins of the text are also in \ili{Hebrew}. 

\ili{Biblical Hebrew} is also present in Magriso's commentary, in the numerous biblical references he gives to support his remarks, a common approach in rabbinic prose.

But \ili{Hebrew} is present far beyond that. In this section, I will analyze borrowings, cases of code-switching and the use of \textit{Ladino}, i.e., expressions and phrases taken from this liturgical calque language designed to help medieval Jewish students understand \ili{Hebrew} syntax while using a Hispanic lexicon (\citealt[71]{bossong1990}, \citealt[48--49]{sephiha1973}).

Before we enter the analysis of loanwords, code-switching, and calque translations, let us take a look at how it is graphically presented in the corpus. 

\ea \label{ch3.7.ex.1}
es ke eskapa de ʾisur de \textbf{GEZEL} porke puede ser ke se venga a yerar enel din i es \textbf{M\textsuperscript{E}ḤAYEV} a el \textbf{ZAKAY} i se yama ke le rova de reʾuven i le da a šimʿon. \\
`In other words, he escapes the prohibition on theft, because it is possible that he could be mistaken in his judgement and accuse the innocent, and this is called stealing from Reuben to give to Simon.'
\z


In (\ref{ch3.7.ex.1}), we can see how integrated \ili{Hebrew} loanwords and code-switching are intertwined with the rest of the text in \ili{Judeo-Spanish}. The words in lowercase are integrated loanwords. The three code-switching cases (in uppercase bold letters) are terms from Jewish law that the author does not explain here, but he does explain them in other parts of the corpus. In this case, the repetition of unknown \ili{Hebrew} terms functions as a pedagogical tool. 


\subsection{Borrowings from Hebrew}\label{ch3.7.ss.3.1}

When we speak of borrowing in \ili{Judeo-Spanish}, we need to distinguish between borrowings from contact languages in the \ili{Ottoman} Empire, such as \ili{Turkish} or \ili{Greek}, relatively few of which are present in our text, and borrowings that can be characterized by various sociolinguistic factors, such as cultural pressure and intensity of contact, among others \citep[37--39]{winford2003}. In the case of \ili{Hebrew}, these are not borrowings that are due to contact between two languages. Rather, it was a traditional, social, and cultural relationship with a language whose texts \ili{Judeo-Spanish} speakers studied, even if not always at a sufficiently high level to allow direct access to the literature. For them, it was a language that was quoted and partially understood but not used in everyday life \citep[63]{schwarzwald_rodrigue1981}. 

Furthermore, Marie-Christine \citet[125]{bornes-varol2003} points out that cases of borrowing in \ili{Judeo-Spanish} can often be traced back to forms that already existed in medieval Spain. According to her, these borrowings are linked to ancient forms that have been abandoned and to certain potentialities of the language that speakers come to reactivate or update. 


The first difficult question is, therefore, how to differentiate borrowings from code-switching. If we consider what \citet[73]{sankoff_poplack_vanniarajan1990} and \citet[1177--1178]{poplack_sankoff1988} say, the vast majority of borrowings eventually become morphologically and syntactically integrated into the host language. These words will occupy the same syntactic positions as the corresponding words in the host language.  


A second main and difficult question is to decide which loans are integrated and which are not. 

One of the tools I used to establish which terms were known and used in \ili{Judeo-Spanish} was Joseph \citeauthor{nehama1997}'s dictionary (\citeyear{nehama1997}). Although this is Salonica \ili{Judeo-Spanish} from the 20th century, whereas Magriso wrote his commentary in Constantinople in the 18th, we find numerous entries in the dictionary that correspond to \ili{Hebrew} terms integrated into \ili{Judeo-Spanish} that I also found in the corpus. 

I cross-referenced all the terms I found with David \citegen{bunis1993} lexicon of \ili{Hebrew} and \ili{Aramaic} elements in \ili{Judeo-Spanish}. In this work I was able to find many elements. The presence of a term in the latter work does not necessarily mean that it is integrated, but rather that it is attested in \ili{Judeo-Spanish} literature. I have concluded that if a term is attested, and appears in Nehama's dictionary, there is a good chance that it was also commonly used in Magriso's time, at least for religious purposes.  

Another useful work is Aharon \citegen{maman2019} \textit{Synoptic dictionary of the \ili{Hebrew} component in Jewish languages}. 

These terms were also cross-referenced with the glossary in Angel \citeauthor{berenguer_amador2016}'s (\citeyear{berenguer_amador2016}) thesis on David M. Atias's \textit{La güerta de oro} (1778), a work slightly later than Magriso's \textit{MEʿAM LOʿEZ} but of a secular nature. We can be sure that the terms that appear in this work were familiar to readers of the \textit{MEʿAM LOʿEZ} at the time. 

This methodology has its limits. It is quite impossible to be sure if a \ili{Hebrew} term was known by \ili{Judeo-Spanish} speakers in the \ili{Ottoman} Empire in the 18th century. The only thing we can be sure of is that if a term appears in all the above-mentioned sources, we know it was attested in \ili{Judeo-Spanish} literature. We can also be sure that rabbis were probably familiar with these terms, especially the religious ones, since they knew \ili{Hebrew}. But the common reader of \textit{MEʿAM LOʿEZ} was probably familiar with a much smaller number of \ili{Hebrew} terms, probably only the most common ones related to simple religious practices. 

In most cases, nouns are the most borrowed elements. Next come verbs and adjectives, then adverbs. This is a common phenomenon in language contact in general \citep[77]{sankoff_poplack_vanniarajan1990}, and is also true for \ili{Jewish languages}, as \citet[178--180]{bar-asher1978} points out regarding Moroccan Jewish Arabic\il{Arabic!Moroccan Jewish}.

Among the nouns (in which I also include constructed states that function as a single term), we have terms such as: \textit{ḥorban} `destruction', \textit{ḥol hamoʿed} `intermediate days between the first and last day of Sukkot or Pesach', \textit{hefqer} `lawless', \textit{haqdamaʰ} `introduction', \textit{haknasat ʾorḥim} `hospitality'.  

As for verbs, we have for \textit{mavṭiaḥ} `to promise', \textit{mavḥin} `to observe'. Often the verbs are part of an expression, as in \textit{maḥṭiʾ ʾet harabim} `induces others to sin', \textit{megaleʰ panim batoraʰ} `interprets the Torah'. Since the \ili{Hebrew} present participle\footnote{\textit{beynoni} \texthebrew{בינוני} in \ili{Hebrew}, which can be translated as `intermediary'.} can also be a noun, we need to analyze case-by-case to know whether it is, in the text, a verb or not. Later on, I will talk about \ili{Hebrew} verbs morphologically integrated into \ili{Judeo-Spanish}.

Adjectives: \textit{m\textsuperscript{e}foraš} `explicit', \textit{ʾasur} `forbidden', \textit{ḥarif} `intelligent' (common in \ili{Judeo-Spanish}). 

Adverbs: \textit{ʾafilu} `same' (common in \ili{Judeo-Spanish}), \textit{bif\textsuperscript{e}raṭ} `particularly' (often appears preceded by the \ili{Hebrew} conjunction \textit{ubif\textsuperscript{e}raṭ}), \textit{bi\b{k}lal} `among others', \textit{dawqaʾ} `only, justly'. 

Certain terms of a liturgical nature are frequent and well-known, especially in the \ili{Judeo-Spanish} world \citep[36]{schwarzwald_rodrigue1987}, which leads us to focus on the more pointed terms, where Magriso's rabbinic prose seems to take precedence over his concern for the reading clarity of his public. 

Among these integrated terms we will find in abundance \textit{ʿaśeret yemey tešuvaʰ} `ten days of repentance',\footnote{The ten days between New Year's Day (Rosh Hashana) and the Day of Repentance (Yom Kippur). Appears in the introduction and 5:3.} \textit{sefer} `book', \textit{kohen gadol} `high priest', the names of the various prayers of the day, and many other terms which are therefore part of the synagogal liturgical world.  

The term \textit{sefer} appears in the corpus sometimes as an integrated loanword in Magriso's \ili{Judeo-Spanish} commentary, sometimes as a code-switching case, in quotations made directly in \ili{Hebrew}. Although in \ili{Judeo-Spanish} the term \textit{sefer} alludes to the Bible, in my corpus Magriso systematically uses the construct state \textit{sefer torah} `book of Torah' (in the introduction, 1:1, 1:15, 4:8, 6:5) except in 6:8 where he writes \textit{sefer tehilim} `the book of Psalms'.

Let us look at the examples: first, the occurrences of \textit{sefer} as a code switching in biblical quotations in \ili{Hebrew}:

\ea \label{ch3.7.ex.2}
\textit{{``}ʾAZ NIDB\textsuperscript{e}RU YRʾEY H' ʾIŠ ʾEL REʿEHU WAYAQŠEV H' WAYIŠMAʿ WAYIKATEV B\textsuperscript{e}SEFER ZIKARON L\textsuperscript{e}FANAW L\textsuperscript{e}YRʾEY H UL\textsuperscript{e}ḤOŠ\textsuperscript{e}VEY Š\textsuperscript{e}MO''} \\
`However, the worshippers of the Lord exhorted one another; the Lord listened and heard, and a record of remembrance was set up before him in favor of those who fear the Lord and respect his name.' (3:3, Malachi 3:16)
\z

\ea \label{ch3.7.ex.3}
{[...]} \textit{eskriven sus palavras enel ``SEFER HAZIKARON'' asigún dize el pasuq ayí: ``ʾAZ NIDB\textsuperscript{e}RU YIRʾEY H ʾIŠ ʾEL REʿEHU WAYIKATEV B\textsuperscript{e}SEFER ZIKARON L\textsuperscript{e}FANAW''}

`{[...]} They write their words in the register of remembrance according to what the verse says: `Nevertheless the worshippers of the Lord exhorted one another; the Lord listened and heard, and a register of remembrance was set up before him.'' (3:7, Malachi 3:16) 
\z

In example (\ref{ch3.7.ex.3}) we see that Magriso uses the expression \textit{sefer hazikaron} `register of remembrance', in reference to the verse he quoted above in the first example (3:3) and then repeats. 

In the following examples, Magriso uses the term \textit{sefer} to refer to the Bible, but adds the term \textit{torah} or \textit{tehilim} when referring to the Psalms. Perhaps this is because Magriso knew \ili{Hebrew}, so for him \textit{sefer} was a common noun, a `book', requiring him to specify which book he was referring to each time.

\ea \label{ch3.7.ex.4}
{[...]} \textit{we\b{k}en al zeʰ hadere\b{k} es en todas las mitswot dela Ley ke enel sefer toraʰ se eskrivyeron en korto i en sifras} {[...]} 

`{[...]} in the same way it was the case in all the commandments of the Law which appear in the book of Torah in a brief and encoded manner {[...]}' (1:1)
\z

\newpage
\ea \label{ch3.7.ex.5}
3 \textit{es de katar kavod a el sefer toraʰ de azerle hey\b{k}al apropyado en lugar alto i no eskupir delantre de él} {[...]} \\
3 `It is to respect the book of Torah, to build a special cabinet for it in a high place and not to spit before it {[...]}' (4:8)
\z

\ea\label{ch3.7.ex.6}
\textit{I ansí topamos en dawid hamele\b{k} ke disho en su sefer tehilim ``ṬOV LI TORAT PI\b{K}A MEʾALFEY ZAHAV WAKASEF''}\\
`Thus we find in [the words of] King David who said in the book of Psalms, `More precious to me is the teaching of your mouth than heaps of gold and silver coins''. (6:8)
\z

As far as morphologically integrated borrowings are concerned, some \ili{Hebrew} verbs appear in the text with a \ili{Judeo-Spanish} verbal ending: 

\textit{kafrar} `to blaspheme, to transgress', from the \ili{Hebrew} root \textit{k-p-r} \texthebrew{כפ''ר}, appears in the infinitive form, but also conjugated in the imperfect subjunctive mood \textit{kafrara} and the simple future tense \textit{kafraré}:

\ea\label{ch3.7.ex.n1}
{kafraré} {ḥas w\textsuperscript{e}šalom} {enel ʿiqar} \\ 
`I will transgress, God forbid, in that which is most important' (4:11)
\z

Interestingly, in this sentence, only the preposition \textit{en} and the definite article \textit{el} are in \ili{Judeo-Spanish}. The simple future verb, \textit{kafraré}, has a \ili{Hebrew} root, with a \ili{Judeo-Spanish} simple future suffix \textit{é}, and then the common expression, \textit{ḥas w\textsuperscript{e}šalom}, and the noun \textit{ʿiqar}, are in \ili{Hebrew}.

\begin{description}
 \item[\textit{darsar}] `to preach', from the \ili{Hebrew} root \textit{d-r-š} \texthebrew{דר''ש}, appears in the infinitive form, as well as in the subjunctive \textit{darse} and the present tense \textit{darsas}. This term appears in the corpus with an [s] rather than a [ʃ] which could mean that the borrowing comes rather from \ili{Arabic} where the root is \textit{d-r-s} \textarab{.درس}

\item[{badkar}] `to examine, to explore', from the \ili{Hebrew} root \textit{b-d-q} \texthebrew{בד''ק}, a single occurrence in the infinitive form.

\item[{harvar}] `to destroy, to beat', from the \ili{Hebrew}/\ili{Arabic} root \textit{ḥ-r-b} \texthebrew{חר''ב}, appears in the infinitive form once with prosthetic \textit{a, aharvarlo}, once without \textit{harvar}, as well as in the gerund, \textit{harvando, harvándoles}.
\end{description}

On the other hand, we have integrated and attested terms for which the author meticulously respects \ili{Hebrew} spelling, even if these terms have been morphologically adapted to \ili{Judeo-Spanish}.

The most striking examples are the terms \textit{baʿal habayit}\footnote{\texthebrew{בעל הבית}} `master of the house' or `father of the family', which in \ili{Judeo-Spanish} is found in the form of \textit{balabay} (Nehama) but in the corpus appears exclusively in its \ili{Hebrew} form.

This is also the case for the term \textit{tsaʿar}\footnote{\texthebrew{צער}} `grief', which in \ili{Judeo-Spanish} is found in the form of \textit{sar} (Nehama).

Other examples:\footnote{The \ili{Judeo-Spanish} forms in these examples are given according to Nehama's dictionary.}

\begin{description}
 \item[{ʾarbaʿ w\textsuperscript{e}ʿeśrim}\footnote{\texthebrew{ארבע ועשרים}}] `twenty-four', the 24 books of the Bible, in \ili{Judeo-Spanish} \textit{arbaáveestrim} or \textit{arbavestrím};

\item[\textit{ʿayin haraʿ}\footnote{\texthebrew{עין הרע}}] `the evil eye', in \ili{Judeo-Spanish} \textit{aynará};

\item[\textit{bat qol}\footnote{\texthebrew{בת קול}}] `divine voice', in \ili{Judeo-Spanish} \textit{badkól};

\item[\textit{behemaʰ}\footnote{\texthebrew{בהמה}}] `animal, beast', in \ili{Judeo-Spanish} \textit{beemá} or \textit{bemá};

\item[\textit{tarbut raʿaʰ}\footnote{\texthebrew{תרבות רעה}}] `deviation' in \ili{Judeo-Spanish} \textit{taburrá}.
\end{description}

Other terms, which the author uses with their \ili{Hebrew} form rather than the \ili{Judeo-Spanish} form, show less significant morphological changes, either only the unification of two terms of a construct-state, as in:

\begin{description}

\item[{ben ʾadam}\footnote{\texthebrew{בן אדם}}] `human being', in \ili{Judeo-Spanish} \textit{benadám}, or pronunciation questions, as in:

\item[{derek ʾerets}\footnote{\texthebrew{דרך ארץ}}] `civility', in \ili{Judeo-Spanish} \textit{dereh eres};

\item[\textit{ha\b{k}anaʰ}\footnote{\texthebrew{הכנה}}] `preparation, aptitude', in \ili{Judeo-Spanish} \textit{ahaná};

\item[{hatslaḥaʰ}\footnote{\texthebrew{הצלחה}}] `success', in \ili{Judeo-Spanish} \textit{aslahá}.
\end{description}

The reason for this is undoubtedly identity issues. Throughout the corpus, the author shows a concern for precision in everything concerning the spelling of words of \ili{Hebrew} origin, certainly because of the sacred nature of this language. And while we have seen that some verbs are morphologically integrated, and that there are also \isi{gender} deviations when the author uses a \ili{Judeo-Spanish} article with a \ili{Hebrew} noun, as far as spelling is concerned, there are no visible fluctuations. This is certainly due to what we said in the introduction about the sacred status of the letters of the \ili{Hebrew} \isi{alphabet}.

\subsection{Code-switching}\label{ch3.7.ss.3.2}

Since our corpus is a commentary on a text that is part of the \isi{Mishna}, Magriso's writing is of a rabbinic style, where \ili{Hebrew} terms, whether integrated or not, are an essential part of the discourse. 


I mentioned that the \textit{MEʿAM LOʿEZ} had an educational purpose. \citet[435--438]{harris1999} gives several categories of code-switching levels (between \ili{Judeo-Spanish} and \ili{English} or \ili{Judeo-Spanish} and \ili{Hebrew}) in contemporary \ili{Judeo-Spanish} in the United States and Israel, such as dates, numbers, occupations, places, countries, adjectives, nouns, expressions, etc. We find the same categories in our text, but unlike what \citet[441]{harris1999} asserts, namely that the code-switching cases of her informants were due to their insufficient knowledge of \ili{Judeo-Spanish}, in our case it is due to the strong presence of \ili{Hebrew}, the religious and therefore identity-related character of the text, as well as its pedagogical character.


\subsubsection{Expressions of rabbinic language and Jewish life}\label{ch3.7.ss.3.2.1}

The most frequent cases of code switching in our text are the insertions of religious expressions. 

\begin{description}

 \item[\textit{ʿalaw hašalom}\footnote{\texthebrew{עליו השלום}}]

The expression \textit{ʿalaw hašalom} `peace be upon him', appears 15 times in the text, always in relation to an important figure in Jewish tradition, such as Kings David or Solomon, Abraham or Moses. It also appears once in the plural, \textit{ʿalehem hašalom}, concerning the Sages.
\end{description}

\ea\label{ch3.7.ex.n2}
sigún disho dawid hamele\b{k} {ʿalaw hašalom} \\
`In the words of King David, peace be upon him'
\z

\begin{description}
 \item[\textit{beʿezrat haʾel}\footnote{\texthebrew{בעזרת האל}}]

This expression, which can be translated as `with God's help' or `if God wills', appears seven times in the corpus, often following a reference to a subject Magriso intends to deal with in the future:
\end{description}

\ea\label{ch3.7.ex.n3}
asigún {se avizará} {beʿezrat haʾel} {más adelantre} \\
`as it will be taught, with God's help, below'
\z

In \textit{mišnaʰ} 6:9 this expression appears with another expression, less common in the text:

\ea\label{ch3.7.ex.n4}
\textit{Keryendo dizir ke B\textsuperscript{e}ʿEZRAT HAʾEL LAʿATID LAVOʾ se tyene de arankar la SIṬRAʾ ʾOḤRAʾ del mundo} \\
`Meaning that, with God's help in the future to come, we will have to tear out the Other Side [the dark side] of the world'
\z

\begin{description}
\item[\textit{DIN W\textsuperscript{e}ḤEŠBON}\footnote{\texthebrew{דין וחשבון}}]

This expression, which can be translated as `report, account', appears 6 times in our corpus.
\end{description}

In the following example we have a \ili{Hebrew} expression with the \ili{Hebrew} coordinating conjunction, \textit{w\textsuperscript{e}}, so it can be considered a code switching case:

\ea\label{ch3.7.ex.n5}
{[...]} lo arebiven para dar din weḥešbon de todos los maʿaśim ke izo en este ʿolam {[...]} \\
`{[...]} He is resurrected to give an account of all the deeds he has done in this world {[...]}' (4:27)
\z

This expression appears again in the very text of \textit{PIRQEY ʾAVOT}, in mišnaʰ 3:1: 

\ea\label{ch3.7.ex.n6}
\textit{HISTAKEL BIŠ\textsuperscript{e}LOŠA\textsuperscript{H} D\textsuperscript{e}VARIM W\textsuperscript{e}ʾEYN ʾATAʾ BAʾ LIYDEY ʿAVERA\textsuperscript{H}, DAʿ MEʾAYN BAʾTA MIṬIPA\textsuperscript{H} S\textsuperscript{e}RUḤA\textsuperscript{H} UL\textsuperscript{e}ʾAN ʾATAʾH HOLE\b{K} LIM\textsuperscript{e}QOM ʿAFAR RIMA\textsuperscript{H} W\textsuperscript{e}TOLAʿA\textsuperscript{H} W\textsuperscript{e}LIFNEY MI ʾATAʾ\textsuperscript{H} ʿATID LITEN DIN W\textsuperscript{e}ḤEŠBON W\textsuperscript{e}GOʾ}.\footnote{\texthebrew{הִסְתַּכֵּל בִּשְׁלשָׁה דְבָרִים וְאִי אַתָּה בָא לִידֵי עֲבֵרָה. דַּע מֵאַיִן בָּאתָ, וּלְאָן אַתָּה הוֹלֵךְ, וְלִפְנֵי מִי אַתָּה עָתִיד לִתֵּן דִּין וְחֶשְׁבּוֹן. מֵאַיִן בָּאתָ, מִטִּפָּה סְרוּחָה, וּלְאָן אַתָּה הוֹלֵךְ, לִמְקוֹם עָפָר רִמָּה וְתוֹלֵעָה. וְלִפְנֵי מִי אַתָּה עָתִיד לִתֵּן דִּין וְחֶשְׁבּוֹן}} \\
`Observe these three things, and you will not fall into transgression: know where you come from, where you are going, and before whom you will one day have to give account. From whence thou comest: from a putrid drop; whither thou goest: in place of dust, decay and worms; and before whom thou shalt one day be accountable.'
\z

Here we must take the complete quotation as a case of flagged code-switching case, since Magriso always anticipates his quotations with a discourse marker, as we will see in the next section.  

In addition, this expression appears with the \ili{Judeo-Spanish} conjunction \textit{i}: \textit{din i ḥešbon}. This may be a borrowing of the terms \textit{din} and \textit{ḥešbon}, morphologically integrated by the coordinating conjunction \textit{i}. 

\ea\label{ch3.7.ex.n7}
dar {din i ḥešbon delantre de él} \\
`Reporting to Him [God]'
\z

\subsubsection{Quotes}\label{ch3.7.ss.3.2.2}


Magriso's commentary contains an enormous number of quotations. These include biblical and Talmudic quotations, as well as repetitions of the \textit{PIRQEY ʾAVOT} text.  


Following \citet[1176]{poplack_sankoff1988} definition, most of the quotations in the corpus can be treated as flagged code-switching cases since Magriso almost always introduces them with the discourse markers \textit{asigún dize el pasuq} `as says the verse' or \textit{ainsí dize el pasuq} `so says the verse'. 

\ea\label{ch3.7.ex.n8}
{[...]} ansí dize el pasuq en malʾaki kapítulo 3: ``ʾAZ NIDB\textsuperscript{e}RU YERʾEY H' ʾIŠ ʾEL REʿEHU WAYAQŠEV H' WAYIŠMAʿ WAYIKATEV B\textsuperscript{e}SEFER ZIKARON L\textsuperscript{e}FANAW L\textsuperscript{e}YERʾEY H' UL\textsuperscript{e}ḤOŠ\textsuperscript{e}VEY Š\textsuperscript{e}MO'' {[...]}\footnote{\texthebrew{אָז נִדְבְּרוּ יִרְאֵי יְהוָה, אִישׁ אֶל-רֵעֵהוּ; וַיַּקְשֵׁב יְהוָה, וַיִּשְׁמָע, וַיִּכָּתֵב סֵפֶר זִכָּרוֹן לְפָנָיו לְיִרְאֵי יְהוָה, וּלְחֹשְׁבֵי שְׁמוֹ.}}

`{[...]} so says the verse in Malachi chapter 3: `Nevertheless the worshippers of the LORD exhorted one another; the LORD heard and listened, and a record of remembrance was set up before him for those who fear the LORD and respect his name {[...]}.'' (3:3)
\z

\ea\label{ch3.7.ex.n9}
{[...]} asigún dize el pasuq ``PALGEY MAYIM LEV MELE\b{K} B\textsuperscript{e}YAD H ʾEL KOL ʾAŠER YAḤPOTS YAṬENU'' {[...]}\footnote{\texthebrew{פַּלְגֵי-מַיִם לֶב-מֶלֶךְ, בְּיַד-יְהוָה עַל-כָּל-אֲשֶׁר יַחְפֹּץ יַטֶּנּוּ.}}

`{[...]} Thus says the verse `The king's heart is like a stream in the hand of the Lord; he directs it wherever he will {[...]}.'' (2:3)
\z

When it is not a biblical quotation, Magriso uses other introductory formulas, for example in his commentary on \textit{mišnaʰ} 2:16 Magriso brings a quotation from a poem by Rabbi Menahem di Lonzano, an \ili{Italian} kabbalist of the 16th--17th century: 

\ea\label{ch3.7.ex.n10}
I sovre esto izo un ŠIR el rav rabi menaḥem di lonzano, ke esta es su nota: \\
``ENOŠ KI TAʾAMIN BAʾEL W\textsuperscript{e}TAKIR, HAKI BIR\textsuperscript{e}KAT H' HIʾ TAʿAŠIR. HABIR\textsuperscript{e}KATO TAVOʾEKA B\textsuperscript{e}HAḤASIR ʿAVODATO UV\textsuperscript{e}HEYOT ʾATS L\textsuperscript{e}HAʿAŠIR ``\footnote{\texthebrew{אנוש כי תאמין באל ותכיר, הכי ברכת ה' היא תעשיר. הברכתו תבואך בהחסיר עבודתו ובהיות אץ להעשיר}} \\
`And on this Rabbi Menahem di Lonzano made a poem, which goes like this: `humans believe in God and acknowledge Him, for God's blessing enriches. Will His blessing come to you if you fail in His work and rush to enrich yourself?''
\z

If it is a passage from the \textit{Gemara} or \textit{Mishna} the formula is \textit{dize la mišnaʰ} `the \textit{Mishna} says' or \textit{dize la gemaraʰ} `the \textit{Gemara} says', for example in 1:6 where Magriso quotes tractate \textit{taʿanit} 7a:

\ea\label{ch3.7.ex.n11}
I esto es lo ke dize la gemaraʰ ``HARBE\textsuperscript{H} TORA\textsuperscript{H} LAMAD\textsuperscript{e}TI MERABOTAY UMIḤAVERAY YOTER UMITALMIDAY YOTER MIKULAM''\footnote{\texthebrew{הַרְבֵּה לָמַדְתִּי מֵרַבּוֹתַי, וּמֵחֲבֵירַי יוֹתֵר מֵרַבּוֹתַי, וּמִתַּלְמִידַי יוֹתֵר מִכּוּלָּן.}} \\
`And that's what the Gemara says: `I learned a lot of Torah from my teachers and from my friends even more and from my students more than anything else.'' 
\z

In other cases, especially when Magriso repeats extracts from the text he is commenting on, i.e., \textit{PIRQEY ʾAVOT}, the code-switching cases are direct, without any introduction: 

\ea\label{ch3.7.ex.n12}
Pero te avizo yo ke aun ke tu lo syerves por amistad i no ʿAL M\textsuperscript{e}NAT L\textsuperscript{e}QABEL P\textsuperscript{e}RAS\footnote{\texthebrew{על מנת לקבל פרס}} pero es neʾeman el ŠYt de pagarte tu ŚAKAR {[...]} \\
(lit.) `But I warn you that even if you serve Him out of friendship and not to obtain a reward, the Holy One Blessed be He is faithful to give you your reward {[...]}' (2:18)
\z

The expression \textit{ʿAL M\textsuperscript{e}NAT L\textsuperscript{e}QABEL P\textsuperscript{e}RAS} comes from a \textit{mišnaʰ} commented on earlier. Magriso repeats it here without translation or explanation. 


The following example presents the expression \textit{MIDA\textsuperscript{H} K\textsuperscript{e}NEGED MIDA\textsuperscript{H}} `measure for measure', which has become proverbial in \ili{Judeo-Spanish}:


\ea\label{ch3.7.ex.n13}
Por ke todo pagan de los syelos MIDAH K\textsuperscript{e}NEGED MIDAH\footnote{\texthebrew{מידה כנגד מידה}}, ʾafilu enel echar de la mar después de matado. \\
`For everything is paid for by Heaven measure for measure, even throwing into the sea after [being] killed.' (2:7).
\z

\subsubsection{Prepositions and adverbs}\label{ch3.7.ss.3.2.3}

Another example is the use of \ili{Hebrew} prepositions and adverbs, which Magriso makes extensive use of in his commentary. 

Several prepositions and adverbs are used systematically by the author. This is due to the style of rabbinic prose, and it also lends a more prestigious character to the text, one of \citeauthor{winford2003}'s (\citeyear[38]{winford2003}) hypotheses regarding borrowings. Let us look at three examples: 

\begin{description}
 \item[]\textit{k\textsuperscript{e}neged}] \texthebrew{כנגד} (PREP. `in opposition to, in relation to', 18 occurrences outside quotations)
\end{description}

\ea\label{ch3.7.ex.n14}
{[...]} ya se avizó en otro lugar ke el ḥešbon de las mitswot de la Ley son K\textsuperscript{e}NEGED de los myembros i los nyervos de la presona, onde son dozyentos i kuarenta i ocho myembros K\textsuperscript{e}NEGED de los 248 mitswot ʿaśeʰ i 365 mitswot loʾ taʿaśeʰ son K\textsuperscript{e}NEGED de 365 nyervos ke tyene la presona {[...]} \\
`{[...]} we have already said elsewhere that the total of the commandments of the Law is in relation to the limbs and nerves of the person, it is 248 limbs in relation to the 248 positive commandments and 365 negative commandments in relation to the 365 nerves that the person has {[...]}' (2:1)
\z

\begin{description}

\item[\textit{ʾadrabaʾ}] \texthebrew{אדרבא} (ADV. `on the contrary', 7 occurrences), common in \ili{Judeo-Spanish} and \textit{bif\textsuperscript{e}raṭ} \texthebrew{בפרט} (ADV. `particularly', 4 occurrences), common in \ili{Judeo-Spanish}
\end{description}

\ea\label{ch3.7.ex.n15}
{[...]} su pasyón fue grande BIF\textsuperscript{e}RAṬ asu vejez ke asperava de endjentrar algún ijo i ʾADRABAʾ kalía azer algunos esforsos i en lugar de esto se kortó las karnes de modo ke se le aflakó las fuersas.

`{[...]} his passion was great, especially at the time of his old age, for he expected to beget a child, but instead he had to exert himself and instead he cut his flesh and his strength was reduced.' (5:3)
\z

\subsubsection{The Hebrew present participle}\label{ch3.7.ss.3.2.4}

The formula including an inflected form of the Spanish verb ser ‘to be’ and a Hebrew present participle is very common in the text.

\ea\label{ch3.7.ex.n16}
{pasa sovre} alguna mitswaʰ es POGEM i gasta un myembro o nyervo suyo. \\
`he transgresses a commandment, he damages and wears out a limb or a nerve.' (2:1)
\z

Another example combines the verb \textit{ser} `to be' in the past tense followed by a non-integrated \ili{Hebrew} present participle, the \ili{Judeo-Spanish} article \textit{las} and an integrated \ili{Hebrew} noun:

\ea\label{ch3.7.ex.n17}
rabenu haqadoš ke fue M\textsuperscript{e}SADER las mišnayot \\
`our holy master who compiled the articles' (2:10)
\z

In example (\ref{ch3.7.ex.n17}) f\textit{ue M\textsuperscript{e}SADER las mišnayot} `compiled the articles', is a calque of \ili{Hebrew}, lacking the preposition de before the definite article \textit{las}.

\subsubsection{The Hebrew past participle}\label{ch3.7.ss.3.2.5}

With less frequency than the present participle, Magriso also uses the \ili{Hebrew} past participle. As with the present participle, the past participle is always preceded by the verb \textit{ser}: 

\ea\label{ch3.7.ex.n18}
\textit{ʾAHUV}\footnote{\texthebrew{אהוב}} `loved' (5 occurrences): \\
es ``ʾAHUV'' kere dizir ke sea amado i kerensyado de la djente \\
``\textit{ʾAHUV}' means let him be loved and pampered by people' (6:5)
\z

\ea\label{ch3.7.ex.n19}
\textit{TALUY}\footnote{\texthebrew{תלוי}} `depends on' (2 occurrences): \\
La tresera ke es ``šalom'' está TALUY en los baʿaley dinim \\
`The third, `\textit{šalom}', depends on the litigants.'
\z

\subsection{The subtle presence of Hebrew through \textit{Ladino} or calque language}\label{ch3.7.ss.3.3}

In addition to the numerous code-switching cases and \ili{Hebrew} terms integrated into \ili{Judeo-Spanish} that are very frequent in Magriso's writing, \ili{Hebrew} is also very present in \textit{Ladino}, in the sense of a calque language \citep[45, 49]{sephiha1973}. This is due to the \ili{Judeo-Spanish} tradition of translating \ili{Hebrew} works into calque language, and in particular the Tractate of Avot, which counts some 50 different translations since the first in Ferrara in 1552 \citep[139]{schwarzwald_rodrigue1985}.

This calque translation tradition comes from medieval Jewish education in Spain, where rabbis calqued \ili{Hebrew} texts with a \ili{Spanish} lexicon in order to teach children \ili{Hebrew}. With time, it became a tradition in itself, and even had an influence on the vernacular language, as we will see in the following pages.

\subsubsection{Total or partial calque translations}\label{ch3.7.ss.3.3.1}

Some of the expressions used by Magriso are calques in whole or in part of \ili{Hebrew} expressions. Here I will separate these calques into semantic expressions on the one hand, and verbs + prepositions on the other:

\subsubsubsection{Semantic expressions}\label{ch3.7.ss.3.3.1.1}

The author uses several \ili{Hebrew} expressions in his commentary. First, we have several terms that exist in \ili{Hebrew} only in the plural form and that in \ili{Judeo-Spanish} are used in the same way, creating what \citet[66]{sephiha1973} calls ``syntactic distortions'':

\begin{description}

\item[Expressions with \textit{vidas} `lives':]
The \ili{Hebrew} term \textit{ḥayim}\footnote{\texthebrew{חיים}} `life' or `lives', exists only in the plural form. The author uses the calque term \textit{vidas} `lives', 27 times in the corpus, most of them in relation to life in the Future World or \textit{ʿolam habaʾ}.\footnote{\texthebrew{עולם הבא}} Otherwise, when referring to a person's life or other matters of current life, the author uses the singular \ili{Judeo-Spanish} \textit{vida} (43 occurrences). 
\end{description}


\ea\label{ch3.7.ex.n20}
komprar las vidas de {ʿolam haba} kon eyos \\
`buy the lives of the future world with them' (2:8)
\z

\begin{description}

\item[Expressions with \textit{fases} `faces':]

In \ili{Hebrew}, the term \textit{panim}\footnote{\texthebrew{פנים}} `face', exists only in the plural form. Magriso uses the calque translation \textit{fases} on five occasions:
\end{description}

\ea\label{ch3.7.ex.n21}
averguensa  fases de su ḥaver delantre de los munchos \\
`he humiliates his companion in front of others. (Introduction).'
\z

Magriso copies the \ili{Hebrew} expression \textit{ʿaz panim}\footnote{\texthebrew{עז פנים}} `shameless', which appears in \textit{mišnaʰ} 5:19 -- \textit{ʿaz panim lagehinom}\footnote{\texthebrew{עז פנים לגיהנם}} `the shameless will go to Gehenna'.

The other expression is also found in \textit{mišnaʰ} 5:19:

\ea\label{ch3.7.ex.n22}
verguensozo de fases \\
`(somone who is) reserved'
\z

\begin{description}

\item[Expressions with \textit{aguas} `waters':]

In \ili{Hebrew} \textit{mayim}\footnote{\texthebrew{מים}} `water', is always plural. In his commentary on \textit{mišnaʰ} 2:3 Magriso quotes a verse from the book of Proverbs 21:1: \textit{PALGEY MAYIM LEV MELE\b{K} B\textsuperscript{e}YAD H' ʾEL KOL ʾAŠER YAḤPOTS YAṬENU}\footnote{\texthebrew{פַּלְגֵי-מַיִם לֶב-מֶלֶךְ, בְּיַד-יְהוָה עַל-כָּל-אֲשֶׁר יַחְפֹּץ יַטֶּנּוּ.}} `The heart of the king is like a stream in the hand of the Lord; he directs it wherever he wills.' In his explanation of this verse, he uses the following expression:
\end{description}

\ea\label{ch3.7.ex.n23}
los pelagos de las aguas \\
`torrents of water'
\z

Interestingly, in addition to the calque \textit{aguas} `waters', he uses de term \textit{pelagos} `torrents', which sounds very close to the \ili{Hebrew} term in the verse quoted \textit{palgey}. 

\begin{description}

\item[Expressions with \textit{syelos} `heavens':]

In \ili{Hebrew}, \textit{šamayim}\footnote{\texthebrew{שמיים}} is another term that exists only in the plural form. It is translated to \ili{Judeo-Spanish} with the plural, \textit{syelos}, when it functions as a synonym for God. When the author speaks of the sky, he uses the singular term, \textit{syelo}. 
\end{description}

\ea\label{ch3.7.ex.n24}
para azer	veluntad de tu padre  ke	  en los syelos \\
`to do [the] will of your father who in the skies' (3:15)
\z

The term \textit{syelos} is translated from the \ili{Hebrew} plural \textit{šamayim} and in the expression \textit{ke en los syelos} `who [is] in the skies', \textit{ke} is a calque of the \ili{Hebrew} junctor \textit{še}. In \ili{Hebrew} we would say \textit{ʾavi\b{k}a šebašamayim}. \citet[447--448]{schwarzwald_rodrigue1989b} notes this calque as well as the entire expression we have here \textit{ke en los syelos} as being a calque of the \ili{Hebrew} expression \textit{šebašamayim}\footnote{\texthebrew{שבשמיים}} `who is in Heaven'. This expression also features an elision of the verb ``to be'', \textit{estar}.

Then we have calque expressions, composed with a verb and/or an adjective:

\begin{description}

\item[\textit{Afirmar mitswot}] `to fulfill the commandments', from \ili{Hebrew} \textit{l\textsuperscript{e}qayem mitswot}:\footnote{\texthebrew{לקיים מצוות}}
\end{description}

\ea\label{ch3.7.ex.n25}
afirmar las mitswot  ʿaśeʰ \\
`fulfill  the positive commandments' (1:2)
\z

The \ili{Judeo-Spanish} verb \textit{afirmar} translates the \ili{Hebrew} verb \textit{l\textsuperscript{e}qayem} `to fulfill' \citep[426]{schwarzwald_rodrigue1989b}. This same calque also appears when translating the \ili{Hebrew} term \textit{mitswot} `commandments', into \ili{Judeo-Spanish} \textit{komandos}.

\ea\label{ch3.7.ex.n26}
travajar en la Ley i afirmar los komandos 	 del	ŠYt \\
`study the Law and fulfill the commandments of the Holy One Blessed be He.'
\z


\begin{description}

\item[\textit{Ayegar al lugar}] `to be in the place of', translated from \ili{Hebrew} \textit{l\textsuperscript{e}hagiʿa lim\textsuperscript{e}qom} (lit. `to arrive in the place'):\footnote{\texthebrew{להגיע למקום}}
\end{description}

In \textit{mišnaʰ} 2:5 Magriso makes a calque translation of part of the original \ili{Hebrew} text he is commenting on:

\ea\label{ch3.7.ex.n27}
no lo djuzges a tu ḥaver lekaf ḥovah si no ayegas a su lugar \\
`don't judge your neighbor negatively until you are in his place.'
\z

Magriso explains \textit{mišnaʰ} 2:5 by copying it as it appears in the \ili{Hebrew} version of the text -- \textit{ʾal tadin ʾet ḥaver\b{k}a ʿad šetagiʿa lim\textsuperscript{e}qomo},\footnote{\texthebrew{אל תדין את חברך עד שתגיע למקומו}} `Do not judge your neighbor until you are in his place.'

Generally speaking, the verb \textit{ayegar} appears in the corpus with the meaning attested by Nehama, namely that of `approaching', `coming close', `arriving', `reaching', `touching the goal'. Here Magriso directly translates the \ili{Hebrew} verb used in the \textit{mišnaʰ -- l\textsuperscript{e}hagiʿa} -- which means `to arrive' by the corresponding verb \textit{ayegar} and thus creates this calque expression. 

He once again uses this verb as a calque in \textit{mišnaʰ} 1:4 :

\ea\label{ch3.7.ex.n28}
komo si	ayegara	 T\textsuperscript{e}MIDIN\footnote{\texthebrew{תמידין}} \\
`as if offering the Perpetual Sacrifice'
\z


Here, this verb translates \ili{Hebrew} verb meaning `to sacrifice', \textit{l\textsuperscript{e}haqriv},\footnote{\texthebrew{להקריב}} whose root, \textit{q-r-b} \texthebrew{קר''ב}, also carries the meaning `to bring near', `to be close'. According to \citet[429]{schwarzwald_rodrigue1989b} it is a Ladinism which also means ``to make an offering.'' 


\begin{description}

\item[God's attributes]

In \textit{mišnaʰ} 4:10 we have another calque translated formula that involves the use of God's name and attributes: 
\end{description}

\ea\label{ch3.7.ex.n29}
\textit{Pero lo derecho es ke no te metas a djuzgar sólo porke no djuzga sólo si no el ken es uno, i su nombre uno ke es el ŠYt bendicho él i bendicho su nombre} \\
`But what is correct is that you do not set out to judge alone, for [he] does not judge alone but he who is one and his name one, who is the Holy One Blessed be He, blessed be he and blessed be his name.'
\z

The form \textit{lo derecho} `the right', translates \ili{Hebrew} \textit{hayašar}\footnote{\texthebrew{הישר}} which is used to indicate what is correct. 

Next, the formula \textit{no djuzga sólo si no el ken es uno, i su nombre uno}. The oppositional preposition \textit{si no} (otherwise) translates the \ili{Hebrew} joncteur \textit{ʾelaʾ}.\footnote{\texthebrew{אלא}}

Finally, the entire formula el ken es uno i su nombre uno el ŠYt bendicho él i bendicho su nombre is a calque of a similar formula in \ili{Hebrew}.\footnote{\textit{YHWH ʾeḥad uš\textsuperscript{e}mo ʾeḥad} from Zechariah 14:9: ``The Lord will be king over all the earth; in that day the Lord will be one and unique will be his name.''

\texthebrew{וְהָיָה יְהוָה לְמֶלֶךְ, עַל-כָּל-הָאָרֶץ בַּיּוֹם הַהוּא, יִהְיֶה יְהוָה אֶחָד וּשְׁמוֹ אֶחָד}

The rest of the formula is modeled on the traditional response we utter when we hear God's name: \textit{baru\textsuperscript{k} huʾ uvaru\b{k} š\textsuperscript{e}mo}, \texthebrew{ברוך הוא וברוך שמו} `blessed be He and blessed be His Name'.} 

\ea\label{ch3.7.ex.n30}
el ken es uno i  su nombre uno el ŠYt  bendicho él i bendicho su nombre \\
`He who is one and his name is one, the HOBH, may he be blessed and his name blessed.'
\z

\subsubsubsection{Translation of prepositions and particles}\label{ch3.7.ss.3.3.1.2}

Another way of copying \ili{Hebrew} is to calque \ili{Hebrew} prepositions. 

\begin{description}

\item[Calques of the \ili{Hebrew} preposition \textit{ʿal}\footnote{\texthebrew{על}} `on, above']
\end{description}

Let us look at some uses of the \ili{Hebrew} preposition \textit{ʿal} regularly used by Magriso:

\ea\label{ch3.7.ex.n31}
pasar sovre una de las mitswot \\
`to break one of the commandments'
\z

Calques \textit{laʿavor ʿal}\footnote{\texthebrew{לעבור על}} `to transgress', (lit.) `to pass over/above'.

We find this usage, with this same meaning also in 3:16: 

\ea\label{ch3.7.ex.n32}
no pasar sovre su neder \\
`don't stifle his vow'
\z

Magriso also uses this preposition in the following sentence: 

\ea\label{ch3.7.ex.n33}
eres ovligado de azerte grande sovre eyos \\
`you have an obligation to lead them' (2:6)
\z

Here \textit{grande sovre eyos} `great over them', in the sense of taking responsibility for others, calques \ili{Hebrew} \textit{gadol ʿal}.\footnote{\texthebrew{גדול על}}  
Another expression modeled on the preposition ʿal is found in 2:7: 

\ea\label{ch3.7.ex.n34}
ise alevantó sovre él ke era más forsudo ke su ḥaver i lo ató pyes i manos para degoyarlo \\
`he got up on him, for he was stronger than his friend, and bound him hand and foot to slit his throat'
\z


Magriso translates here the \ili{Hebrew} expression \textit{laqum ʿal}\footnote{\texthebrew{לקום על}} `to stand up against', (lit.) `to stand on', used to introduce violent action. 

\ea\label{ch3.7.ex.n35}
i se alevantó sovre él \\
`he stood up to him'
\z

\begin{description}

\item[The preposition \textit{ʿal} and the adverb \textit{ʾeyn} `there isn't' or `no':]
\end{description}

In \textit{mišnaʰ} 5:16 we have a religious concept in \ili{Hebrew} \textit{M\textsuperscript{e}ZAKE\textsuperscript{H} ʾET HARABIM}\footnote{\texthebrew{מזכה את הרבים}} `one who acquits the public', in the sense of `obtaining privileges for the public'.

\ea\label{ch3.7.ex.n36}
\textit{Todo el ken es ``M\textsuperscript{e}ZAKEH ʾET HARABIM'' ke los aze azer mitswot i maʿaśim ṭovim i los aze tornar en tešuvaʰ i los enkamina en kamino derecho no pekado vyene por su mano}. \\
`Whoever obtains privileges for the public, that he makes them fulfill commandments and do good deeds and he urges them to repent and puts them on the right path, no sin comes from his hand.'
\z

This explanation of mišnaʰ 5:16 imitates the \ili{Hebrew} formula \textit{ʾeyn ḥeṭʾ baʾ ʿal yado}\footnote{\texthebrew{אין חטא בא על ידו}}:

\ea\label{ch3.7.ex.n37}
no pekado vyene por su mano \\
(Heb.) ʾeyn ḥeṭʾ baʾ ʿal yado \\
`he commits no sin'
\z

It is interesting to see how Magriso places the negation \textit{no} before the noun in place of the indefinite determiner \textit{ningún} `none'. 

Here Magriso translates the preposition \textit{ʿal} rather by the preposition \textit{por}, to give the sense that no sin passes through the hands of one who works for others.  

\subsubsection{Apocopated present participles}\label{ch3.7.ss.3.3.2}

\citet[51--53]{sephiha1973} speaks of apocopated present participles as a way of translating the \ili{Hebrew} present participle into \textit{Ladino}. It is especially in the singular form, according to him, that ``verbs suffer from final wear with tonic accentuation to distinguish them from other homograph verb forms''. The fact that this form existed in Aragonese enabled \textit{Ladino} translators to adopt it freely. 

\ea\label{ch3.7.ex.n38} 
Más da de konsejo dito sinyor ``no seas menospresyán a ningún ombre {[...]}'' \\
`This gentleman also advises us: `Despise no one'' (4:3).
\z

The form \textit{menospresyán} is an apocopated form of \textit{menospresyante}. In this example we see very clearly how this form translates the \ili{Hebrew} verb \textit{baz}, which is a so-called ``monosyllabic'' form (\texthebrew{ע"י/ע"ו} in \ili{Hebrew}) of the verbal pattern \textit{qal}. In \ili{Hebrew} it is written \textit{ʾal t\textsuperscript{e}hi baz l\textsuperscript{e}kol ʾadam}.\footnote{\texthebrew{אל תהי בז לכל אדם}} 

In \textit{mišnaʰ} 4:25 we have two more examples of this phenomenon. The apocopated present participles are \textit{komyén} from \textit{komyente} `he who eats' and \textit{bevyén} from \textit{bevyente} `he who drinks.'

\ea\label{ch3.7.ex.n39}
I el ken deprende Ley de los vyejos si es en koza de p\textsuperscript{e}šaṭ asemeja a komyén uvas dulses i si es koza de sekretos dela Ley asemeja a bevyén vino vyejo {[...]} \\
`And he who learns the Law from the old, if he learns the literal meaning of the text then he resembles the one who eats raisins, and if he learns the hidden meaning of the Law then he resembles the one who drinks old wine {[...]}' \\
\z

\subsubsection{\textit{Ladino} plurals and singular collective nouns}\label{ch3.7.ss.3.3.3}

\citet[65--66]{sephiha1973} and \citet[210]{bornes-varol2008} speak of \textit{Ladino} plurals, which have an archaic form and would express an intensive form of prestige. Here are a few examples:

\ea\label{ch3.7.ex.n40}
kada maka i maka \\
`all the strokes' (5:18) 
\z

\ea\label{ch3.7.ex.n41}
kada tsadiq i tsadiq \\
`all righteous men' (4:12) 
\z

\ea\label{ch3.7.ex.n42}
kada ora  i ora \\
`every hour' (4:27) 
\z

\citet[210]{bornes-varol2008}, referring to Sephiha (\textit{passim}), says that ``this form represents a sophisticated way of speaking''. 

\subsubsection{Ellipses of the verb ``to be''}\label{ch3.7.ss.3.3.4}

In \textit{mišnaʰ} 3:15 we find the following sentence where in addition to the verb ``to be'' the definite article \textit{la} is missing: 

\ea\label{ch3.7.ex.n43}
para azer veluntad de tu padre  ke en los syelos \\
`to do the will of your father in Heaven.'
\z

In \textit{mišnaʰ} 3:21 Magriso directly copies part of the \ili{Hebrew} text, which he comments on: 

\ea\label{ch3.7.ex.n44}
I la butika avyerta i  el butikaryi fyán \\
`And the store is open and the merchant gives credit.'
\z

Directly translating the \ili{Hebrew} text, \textit{haḥanut p\textsuperscript{e}tuḥah wehaḥenwani meqif}\footnote{\texthebrew{החנות פתוחה והחנוני מקיף}} Magriso omits the verb `to be' before the adjective \textit{avyerta} as well as before the apocopated present participle \textit{fyán}. 

\subsubsection{Values of the infinitive form}\label{ch3.7.ss.3.3.5}

We also find infinitive forms with the following values:

\subsubsubsection{The absolute infinitive}\label{ch3.7.ss.3.3.5.1}

The absolute infinitive is a form used in \ili{Biblical Hebrew}, often with a sense of imperative \citep[213]{donnez-guez2013}, but also this paronomastic infinitive serves to accentuate the meaning of the verb in question \citet[241]{bornes-varol2008}. 

\citet[250]{sola-sole1983} points out that one of the most remarkable features of the influence of \ili{Semitic} languages on medieval \ili{Spanish} is the repetition of the verbal concept to mark different nuances, and above all conceptual intensity. Among the various forms he notes are constructions of \ili{Hebrew} origin where the abstract infinitive form is followed by a personal verbal form of the same root. 

In my corpus, I have not found any clear examples of calqued absolute infinitive forms, but the following examples could be forms derived from this infinitive. 

\ea\label{ch3.7.ex.n45}
{[...]} \textit{si no el ovligo ke tyenes es de estar meldando toda tu vida i lo ke puedes alkansar de meldar alkansa} {[...]} \\
`{[...]} otherwise the obligation you have is to study all your life, and what you can achieve by studying, achieve it {[...]}.'
(2:18)
\z

%otherwise, your obligation is to study all your life and %whatever you manage to study, study it (2:18)

This example depends on how we read the sentence. If we read it as in my proposed translation, with the second \textit{alkansa} as an imperative, then we could be looking at a construction reminiscent of the absolute infinitive form \textit{alkansar alkansa}, which is a common form in \ili{Judeo-Spanish}. But it is also possible to read this second \textit{alkansa} as the present tense of the verb \textit{alkansar}. In that case, we would not be dealing with a form reminiscent of the absolute infinitive. 

\ea\label{ch3.7.ex.n46}
Dito sinyor vino a akavidar ke se akavide la persona enel meldar {[...]} \\
`This gentleman came to recommend that the person be careful in studying {[...]}.' (4:16)
\z

The repetition of the verb \textit{akavidar} in example (\ref{ch3.7.ex.n46}) can be seen as a desire on Magriso's part to give greater intensity to his statement, which warns readers as they set out to study. 

There are examples with the verb \textit{azer} `to do':

\ea\label{ch3.7.ex.n47}
I todo el ken aze azer pekar alos munchos {[...]} \\
`he who leads others into sin {[...]}' (5:16)
\z

\ea\label{ch3.7.ex.n48}
se esforsava el yetser haraʿ  kuanto más podía para ver si lo azía azer aver QERI {[...]} \\
`The evil inclination strove to cause him night pollution' (5:4)
\z


In examples (\ref{ch3.7.ex.n47}) and (\ref{ch3.7.ex.n48}), repetition of the verb \textit{azer} is not necessary to give the factitive value: \textit{aze pekar} `he pushes to sin', or \textit{azía aver} `pushes to have', would be sufficient. The addition of the infinitive \textit{azer} reinforces the intensity of the author's words.

There are also constructions that combine gerund and personal verb form, a construction which, according to \citet[252]{sola-sole1983}, is also used to express the absolute infinitive from translations of the Vulgate and Masoretic text into medieval \ili{Spanish}:

\ea\label{ch3.7.ex.n49}
{[...]} en anbezándose a meldar se anbezará el dere\b{k} ʾerets {[...]} \\
`{[...]} by learning to study he will learn respect {[...]}' (3:22)
\z

\ea\label{ch3.7.ex.n50}
{[...]} ke en avyendo daʿat  ay todo bueno -- i en mankando el daʿat le manka todo {[...]} \\
`in having knowledge there is only good, but in lacking knowledge he lacks everything' (2:11)
\z

\subsubsubsection{Nominalized infinitive forms}\label{ch3.7.ss.3.3.5.2}

Nominalized infinitive forms exist in Romance languages \citep[244]{berenguer_amador2016}. 

According to \citet[243]{bornes-varol2008} when infinitive forms are nominalized, they are used in the singular form:

\ea \label{ch3.7.ex.n50b}
{[...]} ``\textit{ké z\textsuperscript{e}\b{k}ut tyenen a mereser regmidos?'' te respondo ke en tu sakar el puevlo de MITSRAYIM tyenen de resivir la Ley i este zekut ya les abasta para regmirlos}.

`{[...]} what right do they have to deserve to be delivered? I answer you that due to your bringing them out of Egypt they must receive the Law and this right is sufficient to deliver them.' (5:6)
\z

Let us look at the infinitive \textit{sacar} `to take':

\ea\label{ch3.7.ex.n51}
en tu  sakar el puevlo de MITSRAYIM \\
`the fact that you took us out of Egypt'
\z

\ea\label{ch3.7.ex.n52}
{[...]} \textit{kere dizir ke tu ojo no pare myentes en otra koza si no todo tu mirar de tus ojos sea eneya.}

`{[...]} it means that your eye should not focus on another thing, but that your gaze should be fixed on it.' (5:21)
\z

\ea\label{ch3.7.ex.n53}
todo tu  mirar 	  de tus ojos sea eneya \\
`the whole gaze of your eyes should be put in her'
\z

\ea
\ea\label{ch3.7.ex.n54}
\textit{Kere dizir ke sea akavidado muncho en el servimyento de los ḥa\b{k}amim ke no se menee de su lado de eyos en su echar i en su alevantar} {[...]} \\
`This means that he must be attentive to the service of the Sages, not straying from them either at their going to bed or at their rising {[...]}' (6:5)
\ex\label{ch3.7.ex.n54b}
\textit{en su echar i  en su  alevantar} \\
in his setting and in his rising
\z
\z

Examples (\ref{ch3.7.ex.n52}) and (\ref{ch3.7.ex.n54}) are clearly calque \ili{Hebrew} constructions. The nominal form \textit{sakar} `the taking out', appears to be a calque of the abbreviated \ili{Hebrew} form \textit{tseʾt}\footnote{\texthebrew{צאת}} `exiting'. The wording of example (\ref{ch3.7.ex.n54}) recalls part of the first paragraph of the prayer \textit{šemaʿ yiśraʾel: w\textsuperscript{e}šinan\textsuperscript{e}tam l\textsuperscript{e}vane\b{k}a w\textsuperscript{e}dibarta bam b\textsuperscript{e}šivt\textsuperscript{e}ka b\textsuperscript{e}veyte\b{k}a uv\textsuperscript{e}le\b{k}t\textsuperscript{e}\b{k}a vadere\b{k} uv\textsuperscript{e}šo\b{k}b\textsuperscript{e}\b{k}a uv\textsuperscript{e}qume\b{k}a} `thou shalt inculcate them [these words] in thy children, thou shalt speak of them (constantly), in thy house or on thy journey, when thou liest down and when thou risest up.'\footnote{\texthebrew{וְשִׁנַּנְתָּם לְבָנֶיךָ וְדִבַּרְתָּ בָּם בְּשִׁבְתְּךָ בְּבֵיתֶךָ וּבְלֶכְתְּךָ בַדֶּרֶךְ וּבְשָׁכְבְּךָ וּבְקוּמֶךָ}}

Example (\ref{ch3.7.ex.n54}) is in line with the custom in Romance languages of nominalizing infinitives, as in \ili{Spanish} and \ili{French}.

This nominalization of the infinitive form is very frequent, and we find el \textit{meldar} `study', \textit{el gastar} `waste' or `expense', \textit{el olvidar} `forgetfulness', \textit{el pezar} `sorrow', \textit{el servir} `service'.

\subsubsubsection{The homological infinitive}\label{ch3.7.ss.3.3.5.3}

\citet[245]{bornes-varol2008} identifies pleonastic constructions that associate a conjugated verb with its corresponding noun. According to \citet[250]{sola-sole1983}, paronomasia is of \ili{Arabic} and \ili{Hebrew} influence. 

This construction is very frequent in my corpus:

\ea\label{ch3.7.ex.n55}
\textit{I syendo ansí guádrate de no avlar tales avlas ael ʿani si no eres ovligado de darle tsedaqaʰ}. \\
\textit{avlar  tales avlas} \\
`to talk such talkings \\
Given this, beware of uttering such words in front of the poor, you are rather obliged to give him charity.' (5:12)
\z

\ea\label{ch3.7.ex.n56}
\textit{11 es ke no podía ninguno kavalyar sovre su kavayo ni asentarse en su siya ni servirse kon su vara ni kon su korona ni kon ninguna koza suya}. \\
\textit{kavalyar sovre su kavayo} \\
`riding  on your horse \\
The eleventh was that no one could ride on his horse or sit on his chair or use his staff or crown or anything else of his.' (6:5)
\z

We find other similar constructions, such as: \textit{kaminar en kamino} `to walk on a path' (6:5), \textit{demandar en esto una demanda} `to ask about it a request' (2:7), \textit{deklarar algún deklaro} `to declare a declaration' (3:16), \textit{depedrer ke depyedre el rašaʿ} `the loss that the wicked loses' (5:1). 

\subsubsubsection{Factitive infinitive forms}\label{ch3.7.ss.3.3.5.4}


Many \ili{Judeo-Spanish} infinitives are built with the prefix \textit{a-}. According to \citet[148]{lapesa1991}, this prefix comes from \ili{Arabic} and gives a factitive value to Castilian verbs. \citet[27]{hanssen1966}, referring to the prosthetic \textit{a-}, also makes reference to this: ``Here we note the influence of word separation errors, the analogical propagation of the preposition \textit{a} and \ili{Arabic} models''.  


In my opinion, to this can be added the fact that Magriso sometimes confuses this prosthetic \textit{a-} with the \ili{Judeo-Spanish} preposition \textit{a} `to', and that constructed infinitive forms in \ili{Hebrew} begin with the preposition \textit{l\textsuperscript{e}}, which is attached to the verb and can be translated, among other things, by the preposition \textit{a} in \ili{Judeo-Spanish}. 

It seems highly likely to me that these elements may have contributed to the appearance of \textit{Ladino} calques with factitive values, as in the following verbs: \textit{akayar} `to be silent' and `to silence', \textit{abashar} `to go down' but also `to make someone or something to go down', \textit{aboltar} `to return', \textit{apezgar} `to make heavy', \textit{asaventar} `to make known or become learned', \textit{ayegar} `to arrive' but also `to sacrifice'.

The verbs \textit{abashar} and \textit{ayegar} (which we have seen earlier as a calque form) deserve particular attention. In the corpus, we find the expression \textit{abashar luvya} (lit.) `to bring down rain', calqued on the \ili{Hebrew} \textit{l\textsuperscript{e}horid gešem}.\footnote{\texthebrew{להוריד גשם}} The \ili{Hebrew} verb \textit{l\textsuperscript{e}horid} is a factitive \textit{hifʿil} form which can be translated as `to bring down'. According to \citet[333]{nehama1997} there is a verb equivalent to `to rain down' in \ili{Judeo-Spanish} which is the verb \textit{luvyar} (a rare verb) or \textit{fazer luvya}.  

To sum up, the presence of prosthetic \textit{a-} in medieval \ili{Spanish}, which, according to \citet[27]{hanssen1966}, combined with the fact that Magriso confuses this \textit{a-} with the \ili{Judeo-Spanish} preposition \textit{a} and the fact that \ili{Hebrew} constructed infinitive forms begin with the preposition \textit{le} which -- as stated earlier -- is equivalent to the preposition \textit{a} in \ili{Judeo-Spanish}, may not only have facilitated the creation of the calques mentioned above, but also contribute to the instability of the position of the preposition \textit{a}, with or without crasis, as we see in the cases of \textit{adizir} versus \textit{a dizir, ameldar} versus a \textit{meldar, akortar} versus \textit{a kortar} and so on. 

\subsubsection{Reduplication}\label{ch3.7.ss.3.3.6}

\subsubsubsection{Homological reduplication}\label{ch3.7.ss.3.3.6.1}

\citet[245]{bornes-varol2008} indicates that reduplication without a copula can be attributed jointly to the influence of \ili{Hebrew} and \ili{Turkish}, where this formulation can be used to mark multitude, when a noun is reduplicated, or to give a superlative or an intensive value if an adjective is reduplicated. 

\ea\label{ch3.7.ex.n57}
I por avizo de esto disho el TANAʾ ``M\textsuperscript{e}ʾOD M\textsuperscript{e}ʾOD HEWE Š\textsuperscript{e}FAL RUAḤ ŠETIQWAT ʾENOŠ RIMA\textsuperscript{H}''. Keryendo dizir muncho muncho deves de ser basho de esprito {[...]} \\
`And on this subject the Sage said `M\textsuperscript{e}ʾOD M\textsuperscript{e}ʾOD HEWE Š\textsuperscript{e}FAL RUAḤ ŠETIQWAT ʾENOŠ RIMAH'. Meaning, many many you must be low in spirit {[...]}' (4:4)
\z

Here the adverb \textit{muncho} is reduplicated to copy the formula ``\textit{M\textsuperscript{e}ʾOD M\textsuperscript{e}ʾOD}'' from the \textit{mišnaʰ} Magriso is explaining. Magriso is deliberately showing his calque of the \ili{Hebrew} form. 


\ea\label{ch3.7.ex.n58}
{[...]} \textit{i le avryeron la tripa i fueron kortando pedasos pedasos de akea enshundya} {[...]} \\
`{[...]} and they cut open his guts and cut off bits and bits  of that fat {[...]}' (5:4)
\z

\ea\label{ch3.7.ex.n59}
{[...]} \textit{fulano i fulano resivyeron de fulano} {[...]} \\
`{[...]} so-and-so received from so-and-so {[...]}' (ML1)
\z

\subsubsubsection{Synonymic reduplication}\label{ch3.7.ss.3.3.6.2}

This is the intensive repetition of synonymous nouns, a frequent phenomenon in \textit{Ladino} Bible translations \citep[246]{bornes-varol2008}. \citet[246--247]{bornes-varol2008}, echoing \citet{sephiha1977}, distinguishes between homosynonymic and heterosynonymic reduplication. 

In my corpus there are abundant examples of both.

\subsubsubsection{Homosynonymic reduplication}\label{ch3.7.ss.3.3.6.2.1}

These are cases where the synonyms come from the same language. They may be nouns or adjectives. 

There are so many instances of this phenomenon, I will just give a few examples from the introduction and the first chapter of the corpus.  

\begin{itemize}
\item \textit{repozo i olgansa} `rest and relaxation' 
\item \textit{santos i temerozos} `saints and fearsome'
\item \textit{kastigeryos i dotrinos} `lessons and teachings'
\item \textit{a las guertas i a los verdjeles} `vegetable gardens and orchards'
\item \textit{sangrías i beverajes i melezinas} `bloodletting, beverages and medicines'
\item \textit{el savyo i el entendido} `the wise and the knowledgeable'
\item \textit{mitswot i maʿśim ṭovim} `commandments and good deeds'
\item \textit{rašaʿ i BAʿAL ʿAVEROT} `the wicked and (the) sinner'
\end{itemize}

\subsubsubsection{Heterosynonymic reduplication}\label{ch3.7.ss.3.3.6.2.2}

These are combinations of nouns that associate languages to which \ili{Judeo-Spanish} may or may not belong \citep[247]{bornes-varol2008}. Sometimes they are not synonyms, but terms that are related:

\begin{itemize}
\item \textit{kastigerios i dotrinos i MIDOT buenas} `lessons and teachings and good attributes' (1:1) (JS + JS + HEB)
\item \textit{ḥaver i amigo} `companion and friend' (1:6) (HEB + JS)
\item \textit{los pashas i el reynado} `kings and royalty' (TURK + JS)
\item \textit{temeles i fuersas} `bases and strengths' or `essentials and strength' (TURK + JS)
\item \textit{los dasyos i las pechas} `taxes and capitation fees' (JS + TURK)
\end{itemize}

\subsubsection{Using the article to highlight the adjective}\label{ch3.7.ss.3.3.7}

As in \ili{Hebrew}, where the adjective following a definite noun is also preceded by a definite article \citep[252]{bornes-varol2008}, we find similar constructions in \ili{Judeo-Spanish} through \textit{Ladino}. However, in my corpus I found only three cases: 

\begin{itemize}
\item \textit{tu ermano el dezeozo} `your brother the eager' (5:12)
\item \textit{tu dezeo el malo} `your desire the evil' (2:4)
\item \textit{su mano la ancha} `his hand the wide one' (3:16)
\end{itemize}

\section{Conclusion}\label{ch3.7.ss.4}

\textit{PIRQEY ʾAVOT}'s commentary in the \textit{MEʿAM LOʿEZ} has a very important heritage value for the \ili{Judeo-Spanish} community. The community used to study the text of \textit{PIRQEY ʾAVOT} during the seven weeks between Passover and the Feast of Weeks (Shavuot), as a kind of means for spiritual elevation. 

As a result, Magriso's commentary in \ili{Judeo-Spanish} makes the text, as well as \textit{PIRQEY ʾAVOT}'s commentaries in the vast Jewish literature, accessible in a simple way to all strata of \ili{Judeo-Spanish} society, including women, which was not a given at the time.  

As we have seen, the use of two \ili{Hebrew} alphabets, the square for the original \ili{Hebrew} text of \textit{PIRQEY ʾAVOT} and the Rashi for the commentary in \ili{Judeo-Spanish} defines -- beyond the fact that it was the norm at the time for this kind of work -- a hierarchy in the structure of the text: the sacred language in square characters and the vernacular language in the ``popular'' \isi{alphabet}. 

Within the commentary, Magriso introduces \ili{Hebrew} far beyond the presence of \ili{Hebrew} in the vernacular, and even in other similar works of the period. For instance, there are 822 different \ili{Hebrew} terms in the text, some of which appear very frequently, not counting the numerous biblical quotations that the author integrates into his commentary. Of these 822 terms, some 345 are not integrated into vernacular \ili{Judeo-Spanish}. 

By way of comparison, in the book \textit{La güerta de oro} published in 1778 in Livorno, we have around 130 \ili{Hebrew} terms. 

In the \textit{MEʿAM LOʿEZ} of Exodus, written partly by Yaakov Huli and partly by Magriso, \ili{Spanish} researcher Aitor García Moreno has identified 567 \ili{Hebrew} terms, in a significantly longer text. 

The high presence of \ili{Hebrew} through a lexicon and calqued structures might raise the question of the pedagogical approach of the author. If he writes a commentary in \ili{Judeo-Spanish} because people do not understand \ili{Hebrew}, then why use so much \ili{Hebrew} in it? The answer to that is a mixture of identity and rabbinical prose. 

\ili{Hebrew} loanwords and code-switching are part of any Jewish language. \ili{Hebrew} has the status of a ``sacred language'', it is the language of the Bible and of the rabbinical literature and is present in the vernacular \ili{Jewish languages} not only through religious life but also through people's Jewish identity. 

If we add to this the fact that rabbinical prose is always made of \ili{Hebrew} terms and references to the Bible, we can explain why a popular commentary as \textit{MEʿAM LOʿEZ} has so much \ili{Hebrew} in it. Possibly the authors thought that in this way people will learn some new \ili{Hebrew} concepts and, above all, the commentary is accessible enough through \ili{Judeo-Spanish}, while keeping its traditional ``Jewish style'' as in all Jewish literature. 

We can add to this that probably because of the ethical content of the \textit{PIRQEY ʾAVOT}, combined with the pedagogical character of the \textit{MEʿAM LOʿEZ}, and probably some personal considerations about the status of the \ili{Hebrew} language, Magriso makes much wider use of \ili{Hebrew} in this commentary.

But, at the same time, Magriso's commentary is less ``reader friendly'' than his predecessor's commentary. Huli's style was also lighter and more moderate. Since Magriso took over Huli's work after the latter passed away, and since they were both disciples of the same rabbi, Judah Rosanès, I can conclude that the difference is due to the different personal styles of the two rabbis, and maybe to the character of the text Magriso commented on; as mentionned above, \textit{PIRQEY ʾAVOT} is not a biblical text, but rather a post-biblical source which deals with ethics. This can also be the reason for Magriso's more moralistic style. 

In this sense, we can presume that using more \ili{Hebrew} words and more rabbinical concepts, sometimes even if there is a risk for readers not to understand it clearly, is part of Magriso's severe commentary style. In addition to that, the fact that he never changes the morphology of integrated loanwords means that \ili{Hebrew} has a special status for him and should have a special place in his commentary. 

This attitude is reflected in Magriso's words, for example, in his commentary for the first article:

\ea\label{ch3.7.ex.n60}
\textit{Ke syendo las palavras de la Ley son muy seradas ke ninguno las entiende, por esto tenemos ovligo de rejirmos asigún loke mos komandan los ḥakamim pork eyos tyenen resivido el deklaro de kada mitswaʰ i mitswaʰ de la Ley}. \\
`The words of the Law [the Bible] are so hermetic that no one can understand them. This is why we are obliged to conduct ourselves according to Wise men's [or rabbis] orders since they have received the meaning of each commandment of the Law.' 
\z

This is the reason for the many borrowings and code-switching cases. We have seen that \ili{Hebrew} is present not only on the lexical level, but also in the morphology of certain verbs, notably apocopated present participles. 

\ili{Hebrew} is also present in semantics, in the many expressions that the author calques, as well as syntactically speaking, in the use of infinitives, prepositions and reduplications. 

At the same time, a question arises which needs further reflection: if, as most of the \ili{Jewish languages} scholars think, \ili{Hebrew} is an inherent element of every Jewish language, could we consider the fact that Magriso chooses not to integrate morphologically \ili{Hebrew} terms, thus giving \ili{Hebrew} a special place of ``sacred language'', to what extent we might speak then of \ili{Hebrew} as a part of the Jewish language? Isn't it, at some point, a disintegration of \ili{Hebrew} from the vernacular language? 

% \section*{Abbreviations}
% \begin{tabularx}{.5\textwidth}{@{}lQ@{}}
% ... & \\
% ... & \\
% \end{tabularx}%
% \begin{tabularx}{.5\textwidth}{@{}lQ@{}}
% ... & \\
% ... & \\
% \end{tabularx}

\sloppy
\printbibliography[heading=subbibliography,notkeyword=this]
\end{document}
